% !TEX root = ../preview.tex

\section{Related Works}
\label{sec:related-works}

% This file mirrors 6_Related_Works.md. Keep both versions aligned.
% Theme 1: Generative modeling and world models.
Diffusion models were initially developed as iterative generative processes that
recover data by reversing a gradual corruption process
\cite{sohl2015deep,ho2020denoising}. Subsequent formulations decoupled the learned
denoiser from a particular stochastic sampler through non-Markovian diffusion and
continuous-time score dynamics \cite{song2021ddim,song2021score}. Flow Matching
provides a closely related continuous generative view: instead of learning the
reverse of a prescribed noising process, it trains a velocity field that
transports a simple source distribution to the data distribution along a chosen
probability path \cite{lipman2023flow}. Rectified Flow further favors increasingly
straight transport trajectories that can be integrated accurately with fewer
solver evaluations \cite{liu2023rectified}. Thus, diffusion and Flow Matching use
different training and path parameterizations, but both generate structured
samples through iterative transport from noise to data.

Their iterative sampling procedures also permit observations to be introduced at
inference time. RePaint conditions an unconditional diffusion prior on known
pixels by modifying only the reverse sampling schedule, without mask-specific
training \cite{lugmayr2022repaint}. More generally, Pseudoinverse-Guided Diffusion
Models ($\Pi$GDM) use a known measurement model and its pseudoinverse to
approximate a conditional score, with a vector--Jacobian product propagating
measurement residuals back to the current sample \cite{song2023pseudoinverse}.
Pokle et al. transfer this correction to pretrained flow models, adding
measurement guidance directly to the velocity field while retaining a
training-free linear inverse solver \cite{pokle2024training}. This progression
establishes the technical basis for using partial measurements to steer a frozen
continuous generator during sampling.

World models shift the generated object from a static sample to the future of an
acting agent. Early neural world models learned compact visual dynamics and used
imagined rollouts for controller training or latent-space planning
\cite{ha2018world,hafner2019planet}. Dreamer and its successors further optimized
policies through trajectories imagined in learned latent dynamics and scaled this
principle across increasingly diverse control domains
\cite{hafner2020dreamer,hafner2021dreamerv2,hafner2025dreamerv3}. With the
emergence of large video generators, UniSim and Genie reframed video prediction
as an interactive, action-conditioned simulation problem, while IRASim introduced
frame-level action conditioning for fine-grained robot--object dynamics
\cite{yang2024unisim,bruce2024genie,zhu2025irasim}. Recent video--action policies
and World-Action Models (WAMs) move one step further by coupling future visual
generation with robot action prediction: GR-2 transfers Internet-video dynamics
into joint video--action fine-tuning, LingBot-VA formulates autoregressive
video--action modeling with closed-loop observation refresh and asynchronous
execution, and DreamZero jointly generates future video and actions from a
pretrained video diffusion backbone
\cite{cheang2024gr2,li2026lingbot,ye2026dreamzero}. Fast-WAM shows that video
co-training can remain beneficial even when explicit future generation is removed
at inference time \cite{yuan2026fastwam}. FBFM addresses the complementary setting
in which such future-state generation is retained, and asks how the generated
state stream can be re-grounded with observations arriving during execution.
